\section{Results}
\label{sec:results}

\subsection{A two-ratio operational model}
\label{sec:regime}

We describe a hierarchical-memory operating point using
\begin{equation}
R_C=\frac{C_{\mathrm{fast}}}{W},\qquad
R_B=\frac{B_{\mathrm{slow}}T_{\mathrm{comp}}}{D}
   =\frac{T_{\mathrm{comp}}}{T_{\mathrm{transfer}}},
\label{eq:ratios}
\end{equation}
where $C_{\mathrm{fast}}$ is usable fast-memory capacity, $W$ is the active
working set, $D$ is compulsory cross-tier traffic per step,
$B_{\mathrm{slow}}$ is sustained bandwidth on the limiting link, and
$T_{\mathrm{transfer}}=D/B_{\mathrm{slow}}$. Unlike a ratio that uses total
step time, $R_B$ is not circular: it compares independently timed computation
and transfer intervals.

We use $\theta_C=0.50$ as an explicit majority-residency convention and derive
$\theta_B=1$ from equality of computation and transfer time. These definitions
produce three operational labels:
\begin{itemize}
\item \textbf{capacity-limited}: $R_C<\theta_C$;
\item \textbf{I/O-limited}: $R_C\geq\theta_C$ and $R_B<\theta_B$; and
\item \textbf{coordination-dominated}: $R_C\geq\theta_C$ and
      $R_B\geq\theta_B$.
\end{itemize}
The capacity boundary is a convention, not a thermodynamic critical point. The
I/O boundary is an overlap boundary: for $R_B<1$, at least
$T_{\mathrm{transfer}}-T_{\mathrm{comp}}$ cannot be hidden by ideal overlap.

\begin{figure}[t]
\centering
\includegraphics[width=0.95\linewidth]{figures/orion_regime_map.png}
\caption{\textbf{Analytical regime map.} The vertical line is the declared
majority-residency convention $\theta_C=0.50$; the horizontal segment is the
derived overlap boundary $\theta_B=1.0$. The diagram is generated from the
same constants as the released classifier and contains no empirical data.}
\label{fig:regime_map}
\end{figure}

\subsection{Irreducible costs and falsifiable predictions}
\label{sec:model}

For accounting purposes, define non-overlapping critical-path contributions so
that a step can be written as
\begin{equation}
T_{\mathrm{total}}=T_{\mathrm{comp}}+T_{\mathrm{mem}}+
T_{\mathrm{swap}}+T_{\mathrm{sync}}.
\label{eq:latency_decomp}
\end{equation}
Independently of that accounting decomposition, a conservative service-time
lower bound that allows the underlying services to overlap is
\begin{equation}
T_{\mathrm{total}}\geq
\max\!\left\{T_{\mathrm{comp}},
\rho W\max(0,1-R_C),\frac{D}{B_{\mathrm{slow}}}\right\},
\label{eq:lower_bound_regime}
\end{equation}
where $W$ denotes bytes accessed in the step and $\rho$ is a lower bound on
service time per non-resident byte for the specified access pattern. Taking the
maximum permits ideal overlap and avoids
double-counting bytes that contribute to both a cache miss and a cross-tier
transfer. Synchronisation and contention can only increase observed time, so
Equation~\eqref{eq:lower_bound_regime} is not a performance predictor. The
residual above it cannot be assigned to one latency component without
event-level measurements.

The formulation makes two falsifiable predictions. First, reducing residency
while holding the access pattern fixed should expose a higher memory-access
cost. Second, when $R_C\geq\theta_C$, configurations with $R_B<1$ must expose
transfer time even under ideal compute--transfer overlap. It does not, without
additional measurements, predict the magnitude of a strategy's speed-up or
claim that strategy rankings necessarily invert.

\subsection{CPU proof-of-concept measurements}
\label{sec:experiments}

We evaluated working-set sensitivity on an Intel Xeon Platinum 8370C CPU
exposed through a virtualised Linux environment. The experiment used a
randomised pointer chain over working sets from 0.5 to 256~MiB and seven trials
per point. A compiled C loop performed dependent loads and timed them with
\texttt{CLOCK\_MONOTONIC\_RAW}; Python was used only to construct the chain and
summarise trials. Linux sysfs reported a 48~MiB last-level cache, which we used
as $C_{\mathrm{fast}}$. Mean dependent-load time was
$\CpuSmallMean\pm\CpuSmallSd$~ns at 0.5~MiB ($R_C=96$),
$\CpuLlcMean\pm\CpuLlcSd$~ns at 48~MiB ($R_C=1$),
$\CpuBoundaryMean\pm\CpuBoundarySd$~ns at 96~MiB ($R_C=0.5$), and
$\CpuLargeMean\pm\CpuLargeSd$~ns at 256~MiB
($R_C=0.188$), where uncertainty is one standard deviation across trials.
The non-monotonic intermediate points show that cache level, translation,
virtualisation, and system noise matter. In particular, these data do not show
a discontinuity at the conventional $R_C=0.5$ label.

As a separate harness check, we used the layered matrix program with 12 layers,
$d=512$, batch size 8, and five timing windows. It varied the fraction of layer
weights retained in the fast tier. The mean total latency averaged over the
three sampled points below $R_C=0.5$ was \MatrixLowMeanMs~ms, compared with
\MatrixFullMeanMs~ms for the fully resident point, a
\MatrixLatencyRatio-fold descriptive ratio. Because computation and copying share
the same CPU memory system, the accompanying $R_B$ sweep was noisy and is not
used as evidence for the I/O boundary.

\begin{table}[t]
\centering
\caption{Selected measured points from the compiled CPU pointer-chain experiment.
Values are mean $\pm$ one standard deviation over seven trials.}
\label{tab:cpu_probe}
\begin{tabular}{rrrr}
\toprule
Working set (MiB) & $R_C$ & Latency (ns/load) & Interpretation\\
\midrule
0.5 & 96.000 & $\CpuSmallMean\pm\CpuSmallSd$ & small working set\\
48  & 1.000  & $\CpuLlcMean\pm\CpuLlcSd$ & reported LLC capacity\\
96  & 0.500  & $\CpuBoundaryMean\pm\CpuBoundarySd$ & residency convention\\
256 & 0.188  & $\CpuLargeMean\pm\CpuLargeSd$ & largest working set\\
\bottomrule
\end{tabular}
\end{table}

These measurements are proof-of-concept results from one CPU, not a substitute
for the previously planned A100, MI250, TPU v4, Inferentia2, and Optane
campaign. Accordingly, we make no quantitative cross-accelerator, classifier-
accuracy, energy, or strategy-inversion claim in this version.

\subsection{Reproducibility boundary}

The repository contains the raw JSONL records, machine-readable summaries, and
scripts used for Table~\ref{tab:cpu_probe}. It also contains a simulated backend
for testing the analysis pipeline. Simulated output is labelled as such and is
not used as empirical evidence. Accelerator support in the measurement harness
is an executable protocol for future validation; it is not evidence that those
accelerators were measured in this study.
